\begin{afterword}
\summaryhead

The introduction to the fourth book, \textsc{Mere Reality}, explains that the earlier books compared humans as goal-seeking systems with evolution and artificial intelligence, and that this book compares humans as physical systems with processes that are not mind-like. We are built of inhuman parts; in Giulio Giorello's words from an interview with Daniel Dennett, we have a soul, but it is ``made of lots of tiny robots.'' It previews the book's sequences and then gives background on two debates the book enters.

The first is consciousness. Better models of a bat's behavior may not tell us what it is like to be a bat, and an unconscious ``zombie'' might behave exactly like us. Nagel and Chalmers argue that third-person models cannot capture first-person experience; Chalmers calls consciousness a ``further fact.'' ``Physicalism 201'' will test the argument.

The second is the measurement problem. Quantum amplitudes evolve deterministically until an observation, when the probabilistic Born rule takes over. Bell asked what qualifies a system to be the ``measurer.'' The Copenhagen school split into several views. Yudkowsky argues for many-worlds; since neither Yudkowsky nor the author of the introduction is a physicist, it lists outside sources for checking the arguments.

\responsehead

The introduction does what an introduction should. It previews the book, sets out two debates in plain terms, and says openly that the physics sequences are written by someone who is not a physicist, with a list of sources for checking them. It does not argue the questions it raises, and it need not; it leaves them to the essays that follow. Most of what it says is accurate: the survey figures in its footnote match the 2009 PhilPapers results, and its account of the measurement problem is the standard one. A few statements about who held which view need qualifying.

On consciousness, the zombies are introduced as automata that ``replicate all the overt behaviors'' of a conscious agent. The zombies of the philosophical argument are physical duplicates, and the \textsc{Stanford Encyclopedia} notes that ``plenty of physicalists'' grant that a merely behavioral duplicate might lack experience. The introduction states the stronger version two paragraphs later, so the reader gets both. The question it then asks, whether to ``ditch physicalism,'' is posed as a question; a reader should know that Nagel himself did not draw that conclusion (``It would be a mistake to conclude that physicalism must be false''), and that the argument against physicalism is Chalmers's.

On quantum mechanics, the view that consciousness causes collapse is usually traced to von Neumann and Wigner. The \textsc{Stanford Encyclopedia} says it is ``usually associated with'' Copenhagen but was ``definitely not Bohr's.'' ``Shut up and calculate'' is David Mermin's summary from 1989, not a slogan of the founders. That many-worlds ``has become very popular in recent decades among physicists'' is given without a source. The small polls available support ``popular'': many-worlds came second in one and third in the other. In each it was held by a minority.

In short: an introduction that sets out two debates fairly and gives readers the means to check the physics sequences; the view that consciousness causes collapse was von Neumann's and Wigner's rather than Bohr's.
\end{afterword}
